\section{Related Work}

\textbf{Semantic Segmentation.}
Semantic segmentation can be seen as an extension of image classification from image level to pixel level. In the deep learning era~\cite{resnet,krizhevsky2012imagenet,simonyan2014very,Li_2019_CVPR,wang2018_mix}, FCN~\cite{fcn} is the fundamental work of semantic segmentation, which is a fully convolution network that performs pixel-to-pixel classification in an end-to-end manner. After that, researchers focused on improving FCN from different aspects such as:  enlarging the receptive field~\cite{pspnet, denseaspp,largekernel,dilate,deeplab,deeplabv2,deeplabv3+}; refining the contextual information~\cite{ocnet,cpn,ocrnet,encnet,zhou2019context,refinenet,contextnet,cgnet,apcnet}; introducing boundary information~\cite{ding2019boundary,bertasius2016semantic,li2020improving,segfix,zhen2020joint,gatescnn,dfn,chen2016semantic}; designing various attention modules~\cite{danet,nonlocal,sanet,ccnet,pan,psanet,emanet,gcnet,trans2seg}; or using AutoML technologies~\cite{squeezenas,fasterseg,dynamicrouting,autodeeplab,nekrasov2019fast}.
These methods significantly improve semantic segmentation performance at the expense of introducing many empirical modules, making the resulting framework computationally demanding and complicated. More recent methods have proved the effectiveness of Transformer-based architectures for semantic segmentation~\cite{setr,trans2seg}. However, these methods are still computationally demanding.


\textbf{Transformer backbones}. ViT~\cite{vit} is the first work to prove that a pure Transformer can achieve state-of-the-art performance in image classification. ViT treats each image as a sequence of tokens and then feeds them to multiple Transformer layers to make the classification.  Subsequently, DeiT~\cite{DETR} further explores a data-efficient training strategy and a distillation approach for ViT. More recent methods such as T2T ViT~\cite{yuan2021tokens}, CPVT~\cite{chu2021conditional}, TNT~\cite{han2021Transformer}, CrossViT~\cite{chen2021crossvit} and LocalViT~\cite{li2021localvit} introduce tailored changes to ViT to further improve image classification performance.

Beyond classification, PVT~\cite{pvt} is the first work to introduce a pyramid structure in Transformer, demonstrating the potential of a pure Transformer backbone compared to CNN counterparts in dense prediction tasks. After that, methods such as Swin~\cite{liu2021swin}, CvT~\cite{wu2021cvt}, CoaT~\cite{xu2021co}, LeViT~\cite{graham2021levit} and Twins~\cite{chu2021twins} enhance the local continuity of features and remove fixed size position embedding to improve the performance of Transformers in dense prediction tasks.


\textbf{Transformers for specific tasks}. DETR~\cite{DETR} is the first work using Transformers to build an end-to-end object detection framework without non-maximum suppression (NMS). Other works have also used Transformers in a variety of tasks such as tracking~\cite{transtrack,trackformer}, super-resolution~\cite{chen2020pre}, ReID~\cite{transreid}, Colorization~\cite{kumar2021colorization}, Retrieval~\cite{el2021training} and multi-modal learning~\cite{clip,unit}. For semantic segmentation, SETR~\cite{setr} adopts ViT~\cite{vit} as a backbone to extract features, achieving impressive performance. However, these Transformer-based methods have very low efficiency and, thus, difficult to deploy in real-time applications. 
